\documentclass[11pt,a4paper]{article}
\usepackage[margin=25mm]{geometry}
\usepackage{amsmath,amssymb,booktabs,siunitx,hyperref,listings,microtype}
\hypersetup{colorlinks=true,linkcolor=blue,urlcolor=blue,citecolor=blue,pdftitle={Report 5 - Friction Model Registry for OpenHPL/OpenHPLjl},pdfauthor={Madhusudhan Pandey}}
\lstdefinelanguage{Julia}{keywords={using,begin,end,function,return,const,module,export,include},sensitive=true,comment=[l]\#,morestring=[b]"}
\lstset{language=Julia,basicstyle=\ttfamily\small,breaklines=true,frame=single,columns=fullflexible,showstringspaces=false}
\title{Report 5 - Friction Model Registry for OpenHPL/OpenHPLjl}
\author{Madhusudhan Pandey}
\date{19 September 2026}
\begin{document}
\maketitle
\begin{abstract}
This report separates hydraulic friction from the RigidPipe momentum equation and introduces a build-time friction-model registry for OpenHPLjl. The objective is to keep the waterway equations reusable while allowing simple quadratic losses, Darcy-Weisbach models, laminar friction, and explicit turbulent correlations to be selected without changing the component architecture.
\end{abstract}

\section{Updated Context}
The first rigid-pipe model used the compact hydropower loss law
\[
H_f=RQ|Q|.
\]
This is useful for control-oriented studies, but the coefficient $R$ hides pipe geometry, roughness, viscosity, and flow regime. OpenHPL includes a Darcy-friction function that distinguishes laminar, transition, and turbulent regimes, so OpenHPLjl now separates friction laws into a reusable function family.

\section{Generalized Concept Model}
The waterway momentum equation remains
\[
\frac{L}{gA}\dot Q=H_{\rm in}-H_{\rm out}-H_f(Q).
\]
Only the constitutive relation $H_f(Q)$ changes. The architecture is therefore
\[
\boxed{\text{RigidPipe momentum balance}+\text{selectable friction law}}.
\]
A friction model is selected when the ModelingToolkit system is constructed.

\section{Other Models Available in the Literature}
The first registry contains
\[
\boxed{\text{none}\rightarrow\text{quadratic}\rightarrow\text{Darcy constant-}f\rightarrow\text{laminar}\rightarrow\text{Haaland / Swamee-Jain}}.
\]
The implicit Colebrook-White equation is retained as a residual function for later solution with a nonlinear solver. More detailed models can later include transitional blending, unsteady friction, local fitting losses, and frequency-dependent wall friction.

\section{Mathematics Involved in the OpenHPL Concept/Model}
The Darcy-Weisbach head loss is
\[
\boxed{H_f=f_D\frac{L}{D}\frac{Q|Q|}{2gA^2}}.
\]
The Reynolds number for a circular pipe is
\[
Re=\frac{4|Q|}{\pi D\nu}.
\]
For laminar flow,
\[
\boxed{f_D=\frac{64}{Re}}.
\]
The Haaland explicit turbulent approximation is
\[
\frac{1}{\sqrt{f_D}}
=
-1.8\log_{10}\left[
\left(\frac{\epsilon}{3.7D}\right)^{1.11}
+
\frac{6.9}{Re}
\right].
\]
The Swamee-Jain form is
\[
f_D=
\frac{0.25}{
\left[
\log_{10}\left(
\frac{\epsilon}{3.7D}
+
\frac{5.74}{Re^{0.9}}
\right)
\right]^2
}.
\]
The Colebrook-White relation is implicit in $f_D$ and therefore naturally belongs to a nonlinear-solve layer.

\section{OpenHPL-Specific Modeling Interpretation}
OpenHPL's pipe model is based on the momentum equation with friction and gravity forces, and its Darcy-friction helper distinguishes flow regimes. The OpenHPLjl implementation follows the same separation of physical momentum balance from friction-factor calculation while keeping the first implementation explicit and easy to test.

\section{Implementation in Julia ModelingToolkit / SciML}
The friction family is stored independently from waterways:

\begin{lstlisting}
src/FrictionModels/
|-- FrictionFunctions.jl
|-- FrictionRegistry.jl
|-- README.md
\end{lstlisting}

The registry is

\begin{lstlisting}
const FRICTION_MODEL_REGISTRY = Dict(
    :none => no_friction,
    :quadratic => quadratic_head_loss,
    :darcy_constant => darcy_head_loss,
    :darcy_laminar => darcy_laminar_head_loss,
    :darcy_haaland => darcy_haaland_head_loss,
    :darcy_swamee_jain => darcy_swamee_jain_head_loss,
)
\end{lstlisting}

A pipe selects a law at construction time:

\begin{lstlisting}
@named pipe = RigidPipe(
    friction = :darcy_haaland,
    L = 1000.0,
    diameter = 3.0,
    nu = 1e-6,
    epsilon = 1e-4,
)
\end{lstlisting}

\section{Model Assembly and Connections}
The RigidPipe momentum model is unchanged conceptually:
\[
\dot Q=\frac{gA}{L}
\left(H_{\rm in}-H_{\rm out}-H_f(Q)\right).
\]
Only $H_f$ is supplied by the selected registry function. This makes the same friction family reusable later in penstocks, tailraces, surge-shaft throttles, and fitting models.

\section{Numerical Experiment / Validation}
For a simple quadratic model,
\[
R=0.5,\qquad Q=2~\si{\meter\cubed\per\second},
\]
the expected signed loss is
\[
\boxed{H_f=0.5\times2\times2=2~\si{\meter}}.
\]
For laminar flow at $Re=1000$,
\[
\boxed{f_D=64/1000=0.064}.
\]
The unit tests also verify that Haaland and Swamee-Jain produce physically plausible positive Darcy factors for a turbulent test point.

\section{Expected Outputs}
Useful outputs from friction studies are
\[
Re(t),\qquad f_D(t),\qquad H_f(t),\qquad \Delta H(t),\qquad Q(t).
\]
Comparison plots should show head loss versus flow for each registered model and the resulting change in waterway damping.

\section{Engineering Interpretation}
The compact coefficient $R$ remains valuable for identified or reduced-order hydropower models. Darcy-Weisbach models are more appropriate when geometry, roughness, and viscosity are known. A registry allows these choices to coexist:
\[
\boxed{\text{identified }R\quad\leftrightarrow\quad\text{physics-based Darcy friction}}.
\]
This is useful for FREKI because friction can later be calibrated from measurements without rewriting the waterway model.

\section*{Conclusion}
OpenHPLjl now treats friction as a model family rather than a fixed term. The next step is to compare the registered laws in one RigidPipe simulation and then add transitional and unsteady-friction options where needed.

\begin{thebibliography}{9}
\bibitem{OpenHPLDarcy}
OpenHPL, Darcy friction factor function.
\url{https://build.openmodelica.org/Documentation/OpenHPL.Functions.DarcyFriction.fDarcy.html}

\bibitem{OpenHPLPipe}
OpenHPL, Pipe model.
\url{https://build.openmodelica.org/Documentation/OpenHPL.Waterway.Pipe.html}

\bibitem{OpenHydraulics}
OpenHydraulics wall friction documentation: Swamee-Jain and Colebrook-White forms.
\url{https://build.openmodelica.org/Documentation/OpenHydraulics.Basic.WallFriction.html}
\end{thebibliography}
\end{document}
